% BHU PYQ -- Transcribed & Corrected
% Paper: ENGAE-41: Academic Writing (Mid Term)
% Exam: B.A. / B.Sc. / B.Com. (Hons.) Semester IV (Mid Term) Examination, 2025-26 (NEP)
% Category: Ability Enhancement Course (AEC)

\documentclass[11pt,a4paper]{article}
\usepackage[a4paper,top=2.5cm,bottom=2.5cm,left=2.8cm,right=2.8cm,headheight=14pt]{geometry}
\usepackage{amsmath,amssymb}
\usepackage[T1]{fontenc}
\usepackage[utf8]{inputenc}
\usepackage{lmodern,microtype,fancyhdr,enumitem,hyperref,lastpage,parskip}

\hypersetup{
    colorlinks=true,
    urlcolor=black,
    linkcolor=black,
    pdftitle={ENGAE-41: Academic Writing (Mid Term) -- BHU Sem IV 2025-26 (NEP)}
}

\pagestyle{fancy}
\fancyhf{}
\renewcommand{\headrulewidth}{0.4pt}
\renewcommand{\footrulewidth}{0.4pt}
\fancyhead[L]{\small   \textit{Banaras Hindu University}}
\fancyhead[R]{\small   \textit{ENGAE-41: Academic Writing (Mid Term)}}
\fancyfoot[C]{\small Page  \thepage\ of \pageref{LastPage}}


\newlist{parts}{enumerate}{2}
\setlist[parts,1]{label=(\Alph*),leftmargin=*,itemsep=3pt,topsep=2pt}

\newcommand{\pts}[1]{\hfill   \textit{[#1]}}


\begin{document}


\begin{center}
  {\large\bfseries BANARAS HINDU UNIVERSITY}\\[2pt]
  {
\normalsize B.A. / B.Sc. / B.Com. (Hons.) --- Semester IV Mid Term Examination, 2025-26 (NEP)}\\[6pt]
  \rule{0.8\linewidth}{0.5pt}\\[6pt]
  {\large\bfseries Ability Enhancement Course (AEC)}\\[4pt]
  {\large\bfseries Paper Code: ENGAE-41 : Academic Writing (Mid Term)}\\[4pt]
  {
\normalsize Time: 1 Hour\hfill Maximum Marks: 30}
\end{center}

\rule{\linewidth}{0.4pt}

\smallskip



\noindent   \textit{Instructions: Answer all 30 questions. Each question carries equal marks (1 mark each).}

\bigskip




\noindent   \textbf{Question 1.} What is the main purpose of a Proposal?\pts{1}

\begin{parts}
  \item To tell a story
  \item To describe a past event
  \item To report old research
  \item To convince someone to accept a plan
\end{parts}

\medskip




\noindent   \textbf{Question 2.} What is the main purpose of a review article?\pts{1}

\begin{parts}
  \item To present new lab research
  \item To evaluate a published work
  \item To write personal experience
  \item To give only data
\end{parts}

\medskip




\noindent   \textbf{Question 3.} A good thesis statement should:\pts{1}

\begin{parts}
  \item Be confusing
  \item Clearly state the main argument
  \item Be written only at the end
  \item Include every detail
\end{parts}

\medskip




\noindent   \textbf{Question 4.} Which source is most reliable for academic research?\pts{1}

\begin{parts}
  \item A random blog
  \item A social media post
  \item A long article online
  \item A peer-reviewed journal article
\end{parts}

\medskip




\noindent   \textbf{Question 5.} Extra materials that do not fit in the main text are placed in:\pts{1}

\begin{parts}
  \item Appendices
  \item Title page
  \item Introduction
  \item Conclusion
\end{parts}

\medskip




\noindent   \textbf{Question 6.} Why is audience analysis important before writing a proposal?\pts{1}

\begin{parts}
  \item To choose font style
  \item To decide page number
  \item To know who will read it and what they need
  \item To make it colorful
\end{parts}

\medskip




\noindent   \textbf{Question 7.} Which is NOT a good quality of technical writing?\pts{1}

\begin{parts}
  \item Clear language
  \item Correct technical terms
  \item Proper headings
  \item Casual conversational style
\end{parts}

\medskip




\noindent   \textbf{Question 8.} Which part of a proposal best explains how to solve the problem?\pts{1}

\begin{parts}
  \item Problem statement
  \item Proposed solution/recommendation
  \item Appendix
  \item Title
\end{parts}

\medskip




\noindent   \textbf{Question 9.} A humanities essay mainly uses:\pts{1}

\begin{parts}
  \item Numbers and graphs
  \item Ideas, interpretation and arguments
  \item Only experiments
  \item Fixed scientific format
\end{parts}

\medskip




\noindent   \textbf{Question 10.} The best way to avoid plagiarism is to:\pts{1}

\begin{parts}
  \item Read the source well and write in your own words
  \item Replace words with synonyms only
  \item Copy everything
  \item Avoid all sources
\end{parts}

\medskip




\noindent   \textbf{Question 11.} Which of these is NOT a citation style?\pts{1}

\begin{parts}
  \item APA
  \item MLA
  \item Latin
  \item Chicago
\end{parts}

\medskip




\noindent   \textbf{Question 12.} A good abstract should:\pts{1}

\begin{parts}
  \item Be informal
  \item Correctly summarize the whole paper
  \item Include only questions
  \item Be very long
\end{parts}

\medskip




\noindent   \textbf{Question 13.} Which is NOT an important pre-writing question?\pts{1}

\begin{parts}
  \item Which font should be used?
  \item Who is the audience?
  \item Why is it written?
  \item How much to include?
\end{parts}

\medskip




\noindent   \textbf{Question 14.} A technical article should:\pts{1}

\begin{parts}
  \item Be emotional
  \item Be only personal opinion
  \item Be for entertainment
  \item Present facts, analysis and conclusions
\end{parts}

\medskip




\noindent   \textbf{Question 15.} Why is an outline useful?\pts{1}

\begin{parts}
  \item No need to revise later
  \item It gives a clear plan of the writing
  \item It guarantees success
  \item It adds references
\end{parts}

\medskip




\noindent   \textbf{Question 16.} Why are Journal articles usually valued more than conference papers?\pts{1}

\begin{parts}
  \item They are longer
  \item They go through stronger peer review
  \item They use easy words
  \item They are always online
\end{parts}

\medskip




\noindent   \textbf{Question 17.} A good academic review should:\pts{1}

\begin{parts}
  \item Be very long
  \item Ignore weaknesses
  \item Give balanced criticism
  \item Always praise the work
\end{parts}

\medskip




\noindent   \textbf{Question 18.} What should be included in an abstract?\pts{1}

\begin{parts}
  \item Problem, method, findings and conclusion
  \item Full tables
  \item All references
  \item Full literature review
\end{parts}

\medskip




\noindent   \textbf{Question 19.} A proposal may be rejected if:\pts{1}

\begin{parts}
  \item It has many sections
  \item The reason for the study is weak
  \item It is formal
  \item It is individual work
\end{parts}

\medskip




\noindent   \textbf{Question 20.} Zotero is useful because it:\pts{1}

\begin{parts}
  \item Gives research topics
  \item Writes the paper
  \item Checks arguments
  \item Organizes and formats references
\end{parts}

\medskip




\noindent   \textbf{Question 21.} Google Docs helps group writing because:\pts{1}

\begin{parts}
  \item Many people can write together at the same time
  \item It removes discussion
  \item It writes automatically
  \item It checks every citation
\end{parts}

\medskip




\noindent   \textbf{Question 22.} A journal style guide gives rules about:\pts{1}

\begin{parts}
  \item Reviewer names
  \item Author age
  \item Teacher approval
  \item Formatting, references, symbols and images
\end{parts}

\medskip




\noindent   \textbf{Question 23.} A research article is different from a review article because it:\pts{1}

\begin{parts}
  \item Only critiques
  \item Is only for conferences
  \item Gives original research findings
  \item Has no method section
\end{parts}

\medskip




\noindent   \textbf{Question 24.} Multimedia should be used in academic writing when:\pts{1}

\begin{parts}
  \item It makes the paper long
  \item It improves understanding
  \item It replaces all writing
  \item It is copied from the internet
\end{parts}

\medskip




\noindent   \textbf{Question 25.} Good feedback on writing should be:\pts{1}

\begin{parts}
  \item Only grammar correction
  \item Only praise
  \item Given at the end only
  \item Specific and helpful for revision
\end{parts}

\medskip




\noindent   \textbf{Question 26.} In digital sources, ``authority'' means:\pts{1}

\begin{parts}
  \item Social media popularity
  \item Government approval
  \item Author's expertise and credentials
  \item Length of article
\end{parts}

\medskip




\noindent   \textbf{Question 27.} A reflective journal is used for:\pts{1}

\begin{parts}
  \item Writing personal learning, difficulties and insights
  \item Copying passages
  \item Writing lecture notes
  \item Listing references
\end{parts}

\medskip




\noindent   \textbf{Question 28.} MLA in-text citation usually includes:\pts{1}

\begin{parts}
  \item Full name and year
  \item Author's last name and page number
  \item Title and publisher
  \item Footnote number
\end{parts}

\medskip




\noindent   \textbf{Question 29.} Which is plagiarism?\pts{1}

\begin{parts}
  \item Using common knowledge
  \item Proper quotation
  \item Correct citation
  \item Using someone's idea without credit
\end{parts}

\medskip




\noindent   \textbf{Question 30.} Good writing assessment should:\pts{1}

\begin{parts}
  \item Compare students only
  \item Check both writing process and final product
  \item Focus only on grammar
  \item Give one short comment
\end{parts}

\vfill
\rule{\linewidth}{0.4pt}

\begin{center}\small
\href{https://bhu-exam.vercel.app/}{Click here to visit our website}\\[4pt]
\href{https://me-aryan.vercel.app/}{Click here to visit our contributor}
\end{center}

\end{document}
